\documentclass[10pt,a4paper]{amsart}
\usepackage[margin=19mm]{geometry}
\usepackage{amsmath,amssymb,mathtools}
\usepackage[scr=rsfs,cal=euler]{mathalfa}
\usepackage{xfrac}
\usepackage[unicode,hidelinks,bookmarks=false]{hyperref}
\newcommand{\ie}{i.e.\ }
\newcommand{\cf}{cf.\ }
\newcommand{\resp}{respectively\ }
\newcommand{\Subsection}{\subsection}
\providecommand{\nobreakdash}{\nobreak\hbox{-}}
\setlength{\emergencystretch}{2em}
\multlinegap0pt
\usepackage{bm}
\usepackage{bbm}
\usepackage{dsfont}
\usepackage{xspace}
\usepackage{cancel}
\usepackage{mathtools}
\usepackage{amsthm}
\makeatletter
\@ifundefined{Theorem}{\newtheorem{Theorem}{Theorem}}{}
\@ifundefined{Proposition}{\newtheorem{Proposition}[Theorem]{Proposition}}{}
\makeatother
\newcommand{\E}{\mathbb{E}}
\title[Credit risk for large portfolios of green and brown loans: e]{Credit risk for large portfolios of green and brown loans: extending the ASRF model}
\author{Alessandro Ramponi, Sergio Scarlatti}
\date{Source: arXiv:2506.12510v1 (CC BY 4.0)}
\begin{document}
\maketitle

\noindent\emph{Source and licence.} Credit risk for large portfolios of green and brown loans: extending the ASRF model. Version arXiv:2506.12510v1 (CC BY 4.0), distributed under the Creative Commons Attribution 4.0 International licence, as recorded in the arXiv licence metadata for this exact version and shown on the abstract page https://arxiv.org/abs/2506.12510v1. Full rights notice: licenses/arxiv-2506.12510-CC-BY-4.0.txt.\par\smallskip

\noindent\textit{Editorial scope.} Counted unit: the Proposition whose source label is \texttt{Limit-distrib} in the article (source0.tex lines 322--344, source label \texttt{Limit-distrib}), delivered together with the article's own complete original proof of it (source0.tex lines 345--377, one \texttt{proof} environment of 33 lines). No mathematics is written, repaired or re-derived: statement and proof are the article's own text line for line. Required context retained uncounted: none. Every definition, display and label of those lines is retained unchanged. Omitted from this package and not redistributed: 1--321, 378--1783. Nothing inside a retained excerpt is omitted or altered: every retained body is delivered exactly as the article's lines are written, so no omission receipt of any kind accompanies this package. Citations: the retained excerpts carry no citation command at all, so no bibliography entry of the article is transcribed and no reference is redistributed. Reference adaptation: none was needed and none was made; every reference inside the delivered unit resolves inside it, so no unresolved reference or citation is produced. Printed slips kept literal: no printed word, symbol, hypothesis or number of the counted statement or of its proof was altered. Layout and class: the article's own class is \texttt{article} with packages including amssymb, amsmath, amsfonts, amsthm, amscd, graphicx, textcomp, color, booktabs, multirow, siunitx, caption; the wrapper is the lane's pinned \texttt{amsart} editorial wrapper at 10pt on A4, which restates the theorem-like environments the retained text needs and adds the article's own macro definitions that the retained text needs (\texttt{\textbackslash E}), with extra wrapper packages (bm, bbm, dsfont, xspace, cancel, mathtools). The delivered numbering is the unit's own. Assets: no retained span contains a figure environment, a graphics-inclusion command, a PDF-page inclusion, a table or a TikZ picture; no image file is copied into this package, the package declares no asset dependency and no image file is redistributed with it. Licence: version arXiv:2506.12510v1 of this article is distributed under the Creative Commons Attribution 4.0 International licence, as recorded in the arXiv licence metadata for this exact version and shown on the abstract page https://arxiv.org/abs/2506.12510v1. A full rights notice is delivered at licenses/arxiv-2506.12510-CC-BY-4.0.txt. Characters: every retained excerpt is pure ASCII, so the delivered engine, XeLaTeX with its default Latin Modern text font, reports no missing character.
\begin{Proposition}\label{Limit-distrib}
 {\bf{(i)}} Suppose $u^b_i = \omega_b/N_b$  and $u^g_j = \omega_g/N_g$, for $i=1,\ldots,N_b$ and $j=1,\ldots,N_g$ , then for $N_b$ and $N_g $ sufficiently large
$$
P(L^N({\bf{u}})\leq \ell) \approx P(L^{mix}_{\omega_b,\omega_g}\leq \ell),\; \forall \ell\in (0,1),
$$
with $L^{mix}_{\omega_b,\omega_g}$ given by the weighted sum
\begin{equation} \label{asymp1}
L^{mix}_{\omega_b,\omega_g}=\omega_bL_b(X^b)+\omega_gL_g(X^g),
\end{equation}
where $L_b(\cdot),L_g(\cdot)$ are explicit functions from $\mathbb{R}$ to $[0,1]$ depending on the parameters of the class and $X^b$ , $X^g$ are skew-normal distributed.\\
 {\bf{(ii)}} If $\delta_b=\delta_g\equiv \delta$ then $X^b\sim X^g$ and we have 
\begin{equation} \label{asymp2}
 L^{mix}_{\omega_b,\omega_g}=v_{\omega_b,\omega_g}(X)
\end{equation}
 with $v_{\omega_b,\omega_g}(\cdot)\equiv \omega_bL_b(\cdot)+\omega_gL_g(\cdot)$ and $X\sim SN(0,1,\frac{\delta}{\sqrt{1-\delta^2}})$. Moreover, the  density of $ L^{mix}_{\omega_b,\omega_g}$ admits the semi-analytical representation given by
\begin{equation}\label{dens}
f^{mix}_{\omega_b,\omega_g}(\ell) =\frac{2 N(\gamma x^*(\ell)) N'(x^*(\ell))}{ \sum_{a = b,g}\left(\frac{\omega_a\rho_a}{\sqrt{1-\rho_a^2}}\right)N'\left(\frac{F_{Y^{(a)}}^{-1}(p^{(a)})-\rho_a x^*(\ell)}{\sqrt{1-\rho_a^2}} \right) },
\end{equation}
over the unit interval, with $\gamma \equiv \frac{\delta}{\sqrt{1-\delta^2}}$, where $x^*(\ell)$ is, for each fixed $\ell \in (0,1)$, the unique root of the equation:
$$
\omega_b\!\underbrace{N\left(\frac{F_{Y^{(b)}}^{-1}(p^{(b)}) - \rho_b x} {\sqrt{1-\rho_b^2}}\right)}_{decreasing \;in\; x}\!\!+\omega_g\!\underbrace{N\left(\frac{F_{Y^{(g)}}^{-1}(p^{(g)}) - \rho_g x} {\sqrt{1-\rho_g^2}}\right)}_{decreasing \;in\; x}\!=\ell.
$$
\end{Proposition}

\begin{proof}
 {\bf{(i):}}: By the previous result is sufficient to show that $L^{mix}_{\omega_b,\omega_g}=\E[L^N(\mathbf{(u^b,u^g)}) | X_1, X_2 ]$ for all $N$. We have
$$
\E[L^N(\mathbf{(u^b,u^g)}) | X_1, X_2 ]=\sum_{i=1}^{N_b}u^b_i\E(\mathbb{I}_{D^b_i}|X_1,X_2)+\sum_{j=1}^{N_g}u^g_j\E(\mathbb{I}_{D^g_i}|X_1,X_2)
$$
$$
=\frac{\omega_b}{N_b}\sum_{i=1}^{N_b}P(D^b_i|X_1,X_2)+\frac{\omega_g}{N_g}\sum_{j=1}^{N_g}P({D^g_j}|X_1,X_2)
$$
$$
=\frac{\omega_b}{N_b}\sum_{i=1}^{N_b}P(Y_i^b\leq K^b|X_1,X_2)+\frac{\omega_g}{N_g}\sum_{j=1}^{N_g}P(Y_j^g\leq K^g|X_1,X_2)
$$
$$
={\omega_b}P(Y_1^b\leq K^b|X_1,X_2)+{\omega_g}P(Y_1^g\leq K^g|X_1,X_2),
$$
where the last equality follows from intra-class homogeneity. Since $Z_1^b$ and $Z_1^g$ are independent from $X_1$ and $X_2$ it follows from ${\bf{(bg3)}}$ that
$$
P(Y_1^a\leq K^a|X_1,X_2)=P\biggl(Z_1^a\leq \frac{K^a-\rho_a(\sqrt{1-\delta^2_a}X_1+\delta_aX_2)}{\sqrt{1-\rho^2_a}}\biggr)=N(h^a(X^a)),
$$
for $a\in \{b,g\}$ and where $N(\cdot)$ denotes the distribution of a standard normal , $h^a(x)\equiv \frac{K^a-\rho_a x}{\sqrt{1-\rho^2_a}}$ and $X^a \equiv \sqrt{1-\delta^2_a} X_1+\delta_a X_2 \sim SN(0,1,\frac{\delta_a}{\sqrt{1-\delta_a^2}})$. Henceforth we have
$$
\E[L^N(\mathbf{(u^b,u^g)}) | X_1, X_2 ]={\omega_b}N(h^b(X^b))+{\omega_g}N(h^g(X^g)),
$$
and the proof ends by setting $L_b\equiv N\circ h^b$ and $L_g\equiv N\circ h^g$; \\
 {\bf{(ii):}} Consider the distribution of the r.v. $v_{\omega_b,\omega_g}(X)$, that is 
 $$
 F_v(\ell)=P(v_{\omega_b,\omega_g}(X)\leq \ell)
 $$
 and notice that the function $x\to v_{\omega_b,\omega_g}(x)$ from  $\mathbb{R}$ to $[0,1]$ is strictly decreasing, therefore for each $\ell\in[0,1]$ the equation $v_{\omega_b,\omega_g}(x)=\ell$ admits a unique solution $x^*(\ell)\equiv v_{\omega_b,\omega_g}^{-1}(\ell)$, henceforth $F_v(\ell)=P(X\geq x^*(\ell))=1-P(X\leq x^*(\ell))$. By differentiating this last expression we obtain 
 $$
 \frac{d F_v(\ell)}{d\ell}=-f_X(x^*(\ell))\bigg(\frac{dv_{\omega_b,\omega_g}^{-1}(\ell)}{d\ell}\bigg)=\frac{f_X(x^*(\ell))}{\sum_{a = b,g}\left(\frac{\omega_a\rho_a}{\sqrt{1-\rho_a^2}}\right)N'\left(\frac{F_{Y^{(a)}}^{-1}(p^{(a)})-\rho_a x^*(\ell)}{\sqrt{1-\rho_a^2}} \right) }\;,
 $$
where $f_X$ denotes the density of $X$. Since this density has the form of the numerator of formula (\ref{dens}) the proof is finished.
\end{proof}

\end{document}
